% Chapter 13: Gaussian vectors.
% Translation of Paola Siri's handwritten notes, PDF pages 203-218 (topic "VETTORI GAUSSIANI (O
% V.A. GAUSSIANE MULTIVARIATE)", whose heading is repeated as "VETTORI GAUSSIANI" on p. 208, where a
% new lecture starts): a recap of Gaussian random variables (p. 203), standard Gaussian vectors
% (pp. 204-205), Gaussian vectors and their characteristic function (pp. 205-208), closure under
% affine maps (pp. 208-210), covariance matrices, the density theorem and the degenerate case, with
% three examples (pp. 210-214), the components of a Gaussian vector, the example of Gaussian
% variables that are not jointly Gaussian, and independence (pp. 215-218).

The Gaussian law is the most important law of probability theory, and this chapter extends it to
random vectors. A \emph{Gaussian vector} is an affine image $A\underline X+\underline\mu$ of a
vector $\underline X$ with independent standard Gaussian components. Its law depends only on two
parameters, the mean vector $\underline\mu$ and the covariance matrix $C=AA^T$, as its
characteristic function shows (\cref{ch:10-characteristic-functions}), and every symmetric positive
semidefinite matrix is the covariance matrix of a Gaussian vector. The family is closed under affine
maps. When $C$ is invertible the vector has an explicit density; when $C$ is singular it lives on a
proper affine subspace and has no density. The components of a Gaussian vector are Gaussian, but
Gaussian random variables need not be jointly Gaussian. Finally, and this is what makes Gaussian
vectors so convenient, their components are independent as soon as they are uncorrelated. The
conditional laws of Gaussian vectors are the subject of \cref{ch:14-gaussian-conditioning}.

Throughout, $(\Omega,\mathcal F,\mathbb P)$ is a probability space, vectors are columns, and
$|\underline x|=(\underline x^T\underline x)^{1/2}$ is the Euclidean norm of
$\underline x\in\mathbb R^N$. $\mathcal M_{N\times K}(\mathbb R)$ is the space of real matrices
with $N$ rows and $K$ columns, $\mathcal M_N(\mathbb R)=\mathcal M_{N\times N}(\mathbb R)$, and
$I_N$ is the identity of $\mathcal M_N(\mathbb R)$.

\section{Gaussian random variables: a recap}\label{sec:gv-recap}

Let us recall the Gaussian random variables. One possible way to define them is the following.

\begin{enumerate}[label=\protect\textcircled{\scriptsize\arabic*},leftmargin=2.2em]
  \item One defines $X\sim\mathcal N(0,1)$ through its density function:
    \[
      f_X(x)=\frac1{\sqrt{2\pi}}\,e^{-\frac{x^2}2}\qquad\forall x\in\mathbb R .
    \]
    \begin{itemize}
      \item $\mathbb E[X]=0$, $\operatorname{var}(X)=1$;
      \item $\Phi_X(t)=e^{-t^2/2}$ for every $t\in\mathbb R$ (\cref{exa:cf-gaussian}).
    \end{itemize}
  \item Let $Y=\sigma X+\mu$ with $\sigma,\mu\in\mathbb R$. $Y$ is called a \emph{Gaussian random
    variable with parameters $\mu$ and $\sigma^2$}, and we write $Y\sim\mathcal N(\mu,\sigma^2)$.%
    \footnote{More precisely, $\mathcal N(\mu,\sigma^2)$ is the law of $\sigma X+\mu$, and
    $Y\sim\mathcal N(\mu,\sigma^2)$ means that $Y$ has this law. It depends on $\sigma$ only through
    $\sigma^2$, because $-X$ has the same law as $X$ (its density is even, or
    \cref{exa:cf-normal-symmetric}): $-\sigma X+\mu$ and $\sigma X+\mu$ have the same law. Added to
    the handwritten notes.}
    \begin{itemize}
      \item $\mathbb E[Y]=\sigma\,\mathbb E[X]+\mu=\mu$;
      \item $\operatorname{var}(Y)=\sigma^2\operatorname{var}(X)=\sigma^2$;
      \item $\Phi_Y(t)=\Phi_{\sigma X+\mu}(t)=e^{it\mu-\frac{\sigma^2t^2}2}$ for every
        $t\in\mathbb R$, by \cref{prop:cf-affine} (this is not needed for the definition).
    \end{itemize}
  \item If $Y\sim\mathcal N(\mu,\sigma^2)$ and $Z=aY+b$ with $a,b\in\mathbb R$, then, taking
    $Y=\sigma X+\mu$,
    \[
      Z=aY+b=a(\sigma X+\mu)+b=(a\sigma)X+(a\mu+b)\ \Longrightarrow\ Z\sim\mathcal N(a\mu+b,a^2\sigma^2).
    \]
  \item \begin{itemize}
      \item If $\sigma=0$, then $Y=\mu$: $Y$ is \emph{degenerate} (its law is the Dirac mass
        $\delta_\mu$).
      \item If $\sigma\ne0$, then $Y$ is absolutely continuous, with density
        (\cref{prop:tr-affine-gaussian,rem:rv-affine-1d})
        \begin{equation}\label{eq:gv-density-1d}
          f_Y(y)=f_X\Big(\frac{y-\mu}\sigma\Big)\frac1{|\sigma|}
          =\frac1{\sqrt{2\pi\sigma^2}}\,e^{-\frac{(y-\mu)^2}{2\sigma^2}}\qquad\forall y\in\mathbb R .
        \end{equation}
    \end{itemize}
\end{enumerate}

\emph{Note.} If $\sigma^2>0$, $Y\sim\mathcal N(\mu,\sigma^2)$ is the absolutely continuous random
variable with density \eqref{eq:gv-density-1d}: for $\sigma^2>0$ one may define
$\mathcal N(\mu,\sigma^2)$ by this density directly, without going through the linear
transformation $\sigma X+\mu$.

In particular, a Gaussian random variable is either absolutely continuous (if $\sigma^2>0$) or
almost surely constant (if $\sigma^2=0$).\footnote{This sentence is added to the handwritten notes.}

\section{Standard Gaussian vectors}\label{sec:gv-standard}

We now generalise to the multivariate case.

\textcircled{\scriptsize 1} Let $X_1,\dots,X_K$ be \iid random variables with law
$\mathcal N(0,1)$, and let
\[
  \underline X=\begin{pmatrix}X_1\\\vdots\\X_K\end{pmatrix}.
\]
Since the $X_i$ are independent, we can compute the density of $\underline X$: for every
$\underline x=(x_1,\dots,x_K)^T\in\mathbb R^K$,
\[
  f_{\underline X}(\underline x)\overset{(a)}{=}\prod_{i=1}^Kf_{X_i}(x_i)
  \overset{(b)}{=}\prod_{i=1}^K\frac1{\sqrt{2\pi}}\,e^{-\frac{x_i^2}2}
  =\frac1{(2\pi)^{K/2}}\,e^{-\frac12\sum_{i=1}^Kx_i^2}
  \overset{(c)}{=}\frac1{(2\pi)^{K/2}}\,e^{-\frac12|\underline x|^2},
\]
where (a) holds because the $X_i$ are independent,\footnote{This is
\cref{thm:pm-density-independence} for $K$ random variables, which has the same proof: if
$X_1,\dots,X_K$ are independent and absolutely continuous, then
$\mathbb P_{\underline X}=\mathbb P_{X_1}\otimes\dots\otimes\mathbb P_{X_K}$, and by Tonelli's
theorem the measure with density $\prod_if_{X_i}(x_i)$ agrees with it on the rectangles, hence
everywhere. Conversely, if the density of $\underline X$ factorises as $\prod_if_i(x_i)$, the
$X_i$ are independent. Added to the handwritten notes.} (b) because each $X_i\sim\mathcal N(0,1)$,
and (c) because $\sum_{i=1}^Kx_i^2=\underline x^T\underline x=|\underline x|^2$.

Observe that the \emph{mean vector} of $\underline X$ is
\[
  \begin{pmatrix}\mathbb E[X_1]\\\vdots\\\mathbb E[X_K]\end{pmatrix}
  =\begin{pmatrix}0\\\vdots\\0\end{pmatrix}=\underline 0\in\mathbb R^K,
\]
while the \emph{variance--covariance matrix} (or \emph{covariance matrix}) of $\underline X$ is
\[
  C=(c_{ij})_{i,j=1,\dots,K}\quad\text{with}\quad
  c_{ij}=\operatorname{cov}(X_i,X_j)=\begin{cases}\operatorname{var}(X_i)=1&\text{if }i=j,\\
  0&\text{if }i\ne j,\end{cases}
  \qquad\text{that is, }c_{ij}=\delta_{ij},
\]
since $X_i\indep X_j$ implies $\operatorname{cov}(X_i,X_j)=0$ for $i\ne j$
(\cref{thm:pm-expectation-product}). Hence $C=I_K$, the identity of $\mathcal M_K(\mathbb R)$.
(For any random vector $\underline Y=(Y_1,\dots,Y_N)^T$ with components in $L^2$, the mean vector
is $(\mathbb E[Y_1],\dots,\mathbb E[Y_N])^T$ and the covariance matrix is
$(\operatorname{cov}(Y_i,Y_j))_{i,j=1,\dots,N}$.)

\begin{definition}[Standard Gaussian vector]\label{def:gv-standard}
A random vector $\underline X$ with values in $\mathbb R^K$ with density function
\[
  f_{\underline X}(\underline x)=\frac1{(2\pi)^{K/2}}\,e^{-\frac12|\underline x|^2}\qquad
  \forall\underline x\in\mathbb R^K
\]
is called a \emph{standard Gaussian vector} (or \emph{standard multivariate Gaussian random
variable}) of dimension $K$, and we write
\[
  \underline X\sim\mathcal N_K(\underline 0,I_K).
\]
\end{definition}

The density factorises as $\prod_{i=1}^K\frac1{\sqrt{2\pi}}e^{-x_i^2/2}$, so the components of a
standard Gaussian vector are independent with law $\mathcal N(0,1)$: the vectors of
\cref{def:gv-standard} are exactly those constructed above.\footnote{By the converse in the
previous footnote, and \cref{prop:pm-marginal-density} for the marginal densities. Added to the
handwritten notes.} We compute the characteristic function of $\underline X$: for every
$\underline t=(t_1,\dots,t_K)^T\in\mathbb R^K$,
\begin{equation}\label{eq:gv-cf-standard}
  \Phi_{\underline X}(\underline t)\overset{(\ast)}{=}\prod_{i=1}^K\Phi_{X_i}(t_i)
  =\prod_{i=1}^Ke^{-\frac{t_i^2}2}=e^{-\frac12\sum_{i=1}^Kt_i^2}=e^{-\frac12|\underline t|^2},
\end{equation}
where $(\ast)$ holds because the $X_i$ are independent (\cref{thm:cf-independence-criterion}).

\section{Gaussian vectors}\label{sec:gv-gaussian}

\textcircled{\scriptsize 2} More generally:

\begin{definition}[Gaussian vector]\label{def:gv-gaussian-vector}
Let $\underline X\sim\mathcal N_K(\underline 0,I_K)$, let $A\in\mathcal M_{N\times K}(\mathbb R)$ and
let $\underline\mu\in\mathbb R^N$. The random vector with values in $\mathbb R^N$
\[
  \underline Y=A\underline X+\underline\mu
\]
is called a \emph{Gaussian vector} (or \emph{multivariate Gaussian random variable}) of dimension
$N$, with parameters $\underline\mu$ and $C=AA^T\in\mathcal M_N(\mathbb R)$, and we write
\[
  \underline Y\sim\mathcal N_N(\underline\mu,C).
\]
\end{definition}

($\underline Y$ is a random vector, since $\underline x\mapsto A\underline x+\underline\mu$ is
continuous; and, as in \textcircled{\scriptsize 2} of \cref{sec:gv-recap}, a random vector with
the same law as $A\underline X+\underline\mu$ is also called a Gaussian vector. We shall see in
\cref{prop:gv-characteristic} that this law depends on $A$ only through $C$.)\footnote{This
paragraph is added to the handwritten notes.}

\begin{remark}[The parameters are the mean vector and the covariance matrix]\label{rem:gv-mean-covariance}
$\underline\mu$ is the mean vector of $\underline Y$, and $C$ is its variance--covariance matrix.
Indeed, write
\[
  \underline Y=A\underline X+\underline\mu\qquad\text{with}\quad A=(a_{ij})_{\substack{i=1,\dots,N\\j=1,\dots,K}},\quad
  \underline\mu=\begin{pmatrix}\mu_1\\\vdots\\\mu_N\end{pmatrix}.
\]
For every $i=1,\dots,N$,
\[
  Y_i=\sum_{j=1}^Ka_{ij}X_j+\mu_i ,
\]
which is in $L^2$, as a linear combination of random variables in $L^2$, and
\[
\begin{aligned}
  \mathbb E[Y_i]&=\sum_{j=1}^Ka_{ij}\underbrace{\mathbb E[X_j]}_{=\,0}+\mu_i=\mu_i,\\
  \operatorname{var}(Y_i)&=\operatorname{var}\Big(\sum_{j=1}^Ka_{ij}X_j\Big)
  \overset{(a)}{=}\sum_{j=1}^Ka_{ij}^2\underbrace{\operatorname{var}(X_j)}_{=\,1}
  =\sum_{j=1}^Ka_{ij}^2=(AA^T)_{ii}=c_{ii},
\end{aligned}
\]
where (a) holds because the $X_j$ are independent. More generally, since the covariance is
bilinear,
\[
\begin{aligned}
  \operatorname{cov}(Y_i,Y_j)&=\operatorname{cov}\Big(\sum_{k=1}^Ka_{ik}X_k+\mu_i,\ \sum_{l=1}^Ka_{jl}X_l+\mu_j\Big)
  =\sum_{k,l=1}^Ka_{ik}a_{jl}\underbrace{\operatorname{cov}(X_k,X_l)}_{=\,\delta_{kl}}\\
  &=\sum_{l=1}^Ka_{il}a_{jl}=(AA^T)_{ij}=c_{ij}.
\end{aligned}
\]
\end{remark}

\begin{remark}[Dimension one]\label{rem:gv-one-dimensional}
If $A\in\mathcal M_{1\times K}(\mathbb R)$, that is, $A=(a_1,\dots,a_K)$ is a row vector, and
$\mu\in\mathbb R$, then $Y=A\underline X+\mu$ is a real random variable:
\[
  Y=\sum_{i=1}^Ka_iX_i+\mu\sim\mathcal N(\mu,\sigma^2),\qquad\text{where }\sigma^2=\sum_{i=1}^Ka_i^2
\]
($Y$ is a linear combination of the \iid $X_i\sim\mathcal N(0,1)$, plus a constant).%
\footnote{The handwritten notes state this without justification. By
\cref{prop:gv-characteristic} below, with $N=1$ and $C=AA^T=\sum_ia_i^2=\sigma^2$,
$\Phi_Y(t)=e^{it\mu-\sigma^2t^2/2}$, which is the characteristic function of
$\mathcal N(\mu,\sigma^2)$ (\textcircled{\scriptsize 2} of \cref{sec:gv-recap}); conclude by
\cref{thm:cf-uniqueness}. When all $a_i\ne0$ one can also use \cref{prop:rv-notable-sums}~a)
repeatedly. Added here.} In particular, $\mathcal N_1(\mu,\sigma^2)=\mathcal N(\mu,\sigma^2)$.\footnote{This
sentence is added to the handwritten notes.}
$Y$ can be degenerate: this happens if $\sum_{i=1}^Ka_i^2=0$, that is, if $a_i=0$ for every
$i=1,\dots,K$, and in that case $Y=\mu$.
\end{remark}

We compute the characteristic function of a Gaussian vector.

\begin{proposition}[Characteristic function of a Gaussian vector]\label{prop:gv-characteristic}
Let $\underline Y\sim\mathcal N_N(\underline\mu,C)$, that is, $\underline Y=A\underline X+\underline\mu$
with $A\in\mathcal M_{N\times K}(\mathbb R)$, $C=AA^T$, $\underline\mu\in\mathbb R^N$ and
$\underline X\sim\mathcal N_K(\underline 0,I_K)$. Then
\begin{equation}\label{eq:gv-cf}
  \Phi_{\underline Y}(\underline t)=e^{i\underline t^T\underline\mu-\frac12\underline t^TC\underline t}
  \qquad\forall\underline t\in\mathbb R^N .
\end{equation}
\end{proposition}

\begin{proof}
By \cref{prop:cf-affine} and \eqref{eq:gv-cf-standard},
\[
  \Phi_{\underline Y}(\underline t)=\Phi_{A\underline X+\underline\mu}(\underline t)
  =e^{i\underline t^T\underline\mu}\,\Phi_{\underline X}(A^T\underline t)
  =e^{i\underline t^T\underline\mu}\,e^{-\frac12|A^T\underline t|^2}
  =e^{i\underline t^T\underline\mu-\frac12(A^T\underline t)^T(A^T\underline t)}
  =e^{i\underline t^T\underline\mu-\frac12\underline t^TC\underline t},
\]
because $(A^T\underline t)^T(A^T\underline t)=\underline t^TAA^T\underline t=\underline t^TC\underline t$.
\end{proof}

In particular, by \cref{thm:cf-uniqueness}, the law of $\underline Y$ depends on $A$ (and on $K$)
only through $C=AA^T$: two Gaussian vectors with the same parameters $\underline\mu$ and $C$ have
the same law. This is what justifies the notation $\mathcal N_N(\underline\mu,C)$ for this
law.\footnote{The handwritten notes leave this consequence implicit; it is used in
\cref{rem:gv-existence} and in the proof of \cref{thm:gv-density}. Added here.} For $N=K$,
$\underline\mu=\underline 0$ and $A=I_K$ we recover the standard Gaussian vector, and a vector
$\underline Y$ has the law $\mathcal N_K(\underline 0,I_K)$ of \cref{def:gv-gaussian-vector} if and
only if its characteristic function is $e^{-|\underline t|^2/2}$, that is, if and only if it is a
standard Gaussian vector in the sense of \cref{def:gv-standard}: the two notations agree.

\section{Affine transformations}\label{sec:gv-affine}

\textcircled{\scriptsize 3}

\begin{theorem}[Closure under affine maps]\label{thm:gv-affine}
The family of multivariate Gaussian laws is closed under affine transformations. Precisely: if
$\underline Y\sim\mathcal N_N(\underline\mu,C)$, $B\in\mathcal M_{M\times N}(\mathbb R)$ and
$\underline\gamma\in\mathbb R^M$, then
\[
  \underline U=B\underline Y+\underline\gamma\sim\mathcal N_M(B\underline\mu+\underline\gamma,\,BCB^T).
\]
\end{theorem}

\begin{proof}
Let $\underline Y$ be a Gaussian vector, $\underline Y\sim\mathcal N_N(\underline\mu,C)$, that is,
$\underline Y=A\underline X+\underline\mu$ with $A\in\mathcal M_{N\times K}(\mathbb R)$, $C=AA^T$,
$\underline\mu\in\mathbb R^N$ and $\underline X\sim\mathcal N_K(\underline 0,I_K)$. Let
$B\in\mathcal M_{M\times N}(\mathbb R)$ and $\underline\gamma\in\mathbb R^M$. Then
$\underline U=B\underline Y+\underline\gamma$ is a random vector with values in $\mathbb R^M$, and
\[
  \underline U=B\underline Y+\underline\gamma=B(A\underline X+\underline\mu)+\underline\gamma
  =(BA)\underline X+(B\underline\mu+\underline\gamma),
\]
with $BA\in\mathcal M_{M\times K}(\mathbb R)$ and $B\underline\mu+\underline\gamma\in\mathbb R^M$.
Hence $\underline U$ is a Gaussian vector, $\underline U\sim\mathcal N_M(\underline m,\Sigma)$, where
\[
  \underline m=B\underline\mu+\underline\gamma,\qquad
  \Sigma=(BA)(BA)^T=B\,AA^T\,B^T=BCB^T .
\]
So $\underline U\sim\mathcal N_M(B\underline\mu+\underline\gamma,\,BCB^T)$.\footnote{If
$\underline Y$ only has the law $\mathcal N_N(\underline\mu,C)$, then $\underline U$ has the law of
$B(A\underline X+\underline\mu)+\underline\gamma$, and the conclusion is the same.}
\end{proof}

\begin{example}[An affine image]\label{exa:gv-affine}
Let $\underline X=(X_1,X_2,X_3)^T\sim\mathcal N_3(\underline\mu,C)$ with
\[
  \underline\mu=\underline 0,\qquad
  C=\begin{pmatrix}2&1&-1\\1&2&1\\-1&1&3\end{pmatrix}.\footnotemark
\]
\footnotetext{The handwritten notes take $c_{33}=1$. Then $\det C=-3<0$, so one eigenvalue is
negative: that matrix is not positive semidefinite, so it is not the covariance
matrix of any random vector (\cref{lem:gv-covariance-psd}) and there is no
$\underline X\sim\mathcal N_3(\underline 0,C)$. With $c_{33}=3$ we have $\det C=3$, $C$ is positive
definite, and only $\operatorname{var}(Z_2)$ below changes, from $1$ to $3$. Corrected here; see
also \exref{ex:gv-covariance-matrices}.}
What is the law of
\[
  \underline Z=\begin{pmatrix}X_1-X_2\\X_1+X_3\end{pmatrix}?
\]
We have $\underline Z=A\underline X+\underline b$, where $\underline b=\underline 0$ and
\[
  A=\begin{pmatrix}1&-1&0\\1&0&1\end{pmatrix}\in\mathcal M_{2\times3}(\mathbb R):\qquad
  \begin{pmatrix}X_1-X_2\\X_1+X_3\end{pmatrix}=\begin{pmatrix}1&-1&0\\1&0&1\end{pmatrix}
  \begin{pmatrix}X_1\\X_2\\X_3\end{pmatrix}.
\]
By \cref{thm:gv-affine}, $\underline Z\sim\mathcal N_2(\underline m,\Sigma)$, where
$\underline m=A\underline\mu+\underline b=\underline 0$ and
\[
\begin{aligned}
  \Sigma=ACA^T&=\begin{pmatrix}1&-1&0\\1&0&1\end{pmatrix}
  \begin{pmatrix}2&1&-1\\1&2&1\\-1&1&3\end{pmatrix}
  \begin{pmatrix}1&1\\-1&0\\0&1\end{pmatrix}\\
  &=\begin{pmatrix}1&-1&-2\\1&2&2\end{pmatrix}\begin{pmatrix}1&1\\-1&0\\0&1\end{pmatrix}
  =\begin{pmatrix}2&-1\\-1&3\end{pmatrix}.
\end{aligned}
\]
Alternatively, since $\underline Z$ is a Gaussian vector, its law is determined by its mean vector
and its covariance matrix, which we can compute directly:
\[
\begin{aligned}
  \mathbb E[Z_1]&=\mathbb E[X_1-X_2]=\mathbb E[X_1]-\mathbb E[X_2]=0,\qquad
  \mathbb E[Z_2]=\mathbb E[X_1+X_3]=\mathbb E[X_1]+\mathbb E[X_3]=0,\\
  \operatorname{var}(Z_1)&=\operatorname{var}(X_1-X_2)=\operatorname{var}(X_1)+\operatorname{var}(X_2)-2\operatorname{cov}(X_1,X_2)=2+2-2=2,\\
  \operatorname{var}(Z_2)&=\operatorname{var}(X_1+X_3)=\operatorname{var}(X_1)+\operatorname{var}(X_3)+2\operatorname{cov}(X_1,X_3)=2+3-2=3,\\
  \operatorname{cov}(Z_1,Z_2)&=\operatorname{cov}(X_1-X_2,X_1+X_3)\\
  &=\underbrace{\operatorname{cov}(X_1,X_1)}_{=\,\operatorname{var}(X_1)}-\operatorname{cov}(X_2,X_1)+\operatorname{cov}(X_1,X_3)-\operatorname{cov}(X_2,X_3)
  =2-1-1-1=-1 .
\end{aligned}
\]
\end{example}

\section{The density of a Gaussian vector}\label{sec:gv-density}

\textcircled{\scriptsize 4} \emph{Problem}: is a Gaussian vector absolutely continuous? If so, what
does its density look like?

\begin{lemma}[Covariance matrices]\label{lem:gv-covariance-psd}
The variance--covariance matrix of any random vector $\underline X$ with values in $\mathbb R^N$
(with components in $L^2$) is symmetric and positive semidefinite.
\end{lemma}

(Note that every matrix with these two properties can be regarded as the variance--covariance
matrix of some random vector: see \cref{rem:gv-existence}.)

\begin{proof}
\begin{itemize}
  \item $C=(c_{ij})_{i,j=1,\dots,N}$ with $c_{ij}=\operatorname{cov}(X_i,X_j)=\operatorname{cov}(X_j,X_i)=c_{ji}$,
    so $C$ is symmetric.
  \item Let $\underline a=(a_1,\dots,a_N)^T\in\mathbb R^N$. Then
    \[
    \begin{aligned}
      \underline a^TC\underline a&=\sum_{i,j=1}^Nc_{ij}a_ia_j=\sum_{i,j=1}^N\operatorname{cov}(X_i,X_j)\,a_ia_j\\
      &=\sum_{i=1}^N\operatorname{var}(X_i)\,a_i^2+2\sum_{i<j}\operatorname{cov}(X_i,X_j)\,a_ia_j
      =\operatorname{var}\Big(\sum_{i=1}^Na_iX_i\Big)\ge0 .
    \end{aligned}
    \]
\end{itemize}
(In our case, $C=AA^T$ and $\underline a^TC\underline a=\underline a^TAA^T\underline a
=(A^T\underline a)^T(A^T\underline a)=|A^T\underline a|^2\ge0$.)
\end{proof}

\begin{theorem}[Density of a Gaussian vector]\label{thm:gv-density}
Let $\underline Y\sim\mathcal N_N(\underline\mu,C)$. If $C$ is invertible, then $\underline Y$ is
absolutely continuous, and its density function is
\begin{equation}\label{eq:gv-density}
  f_{\underline Y}(\underline y)=\frac1{(2\pi)^{N/2}\sqrt{\det C}}\,
  \exp\Big[-\frac12(\underline y-\underline\mu)^TC^{-1}(\underline y-\underline\mu)\Big]
  \qquad\forall\underline y\in\mathbb R^N .
\end{equation}
\end{theorem}

(Here ``absolutely continuous'' means that $\mathbb P_{\underline Y}$ is absolutely continuous
with respect to Lebesgue measure on $\mathbb R^N$, as in \cref{def:pm-ac-vector} for $N=2$.)

\begin{proof}
Observe that $C$ is invertible and positive semidefinite (\cref{lem:gv-covariance-psd}), hence
positive definite;\footnote{The eigenvalues of $C$ are $\ge0$, and none is $0$ because $C$ is
invertible.} moreover, $C$ is symmetric. Therefore there exists $A\in\mathcal M_N(\mathbb R)$
symmetric and positive definite (hence invertible) such that
\[
  C=A^2
\]
(the \emph{principal square root} of $C$).\footnote{By the spectral theorem, $C=Q\Lambda Q^T$ with
$Q$ orthogonal and $\Lambda=\operatorname{diag}(\lambda_1,\dots,\lambda_N)$, $\lambda_i>0$; then
$A=Q\Lambda^{1/2}Q^T$, with $\Lambda^{1/2}=\operatorname{diag}(\sqrt{\lambda_1},\dots,\sqrt{\lambda_N})$,
is symmetric, has the eigenvalues $\sqrt{\lambda_i}>0$, and $A^2=Q\Lambda Q^T=C$. Added to the
handwritten notes.} Set
\[
  \underline X=A^{-1}(\underline Y-\underline\mu)=A^{-1}\underline Y-A^{-1}\underline\mu .
\]
By \cref{thm:gv-affine}, $\underline X$ is a Gaussian vector of $\mathbb R^N$:
\[
  \underline X\sim\mathcal N_N\big(A^{-1}\underline\mu-A^{-1}\underline\mu,\ A^{-1}C(A^{-1})^T\big)
  =\mathcal N_N(\underline 0,I_N),
\]
because $(A^{-1})^T=A^{-1}$ ($A$ is symmetric) and $A^{-1}CA^{-1}=A^{-1}AAA^{-1}=I_N$ ($C=A^2$).
So $\underline X$ is a standard Gaussian vector, with the density of \cref{def:gv-standard} (see the
end of \cref{sec:gv-gaussian}). Note that
\[
  \underline Y=A\underline X+\underline\mu\qquad\text{with $A$ invertible.}
\]
Hence $\underline Y$ is absolutely continuous, and its density function is\footnote{This is
\cref{cor:rv-linear} in $\mathbb R^N$, with the same proof: by the change-of-variables formula in
$\mathbb R^N$ (\cref{thm:rv-jacobi} in $\mathbb R^N$; \citet[Theorem~2.47]{folland}) applied to
$\underline x\mapsto A\underline x+\underline\mu$, whose Jacobian determinant is $\det A$,
$\mathbb P(\underline Y\in B)=\int_{\mathbb R^N}\1_B(A\underline x+\underline\mu)f_{\underline X}(\underline x)\,\dd\underline x
=\int_B f_{\underline X}\big(A^{-1}(\underline y-\underline\mu)\big)\frac{\dd\underline y}{|\det A|}$
for every $B\in\mathcal B(\mathbb R^N)$. Added to the handwritten notes.}
\[
\begin{aligned}
  f_{\underline Y}(\underline y)&=\frac1{|\det A|}\,f_{\underline X}\big(A^{-1}(\underline y-\underline\mu)\big)
  \overset{(a)}{=}\frac1{|\det A|}\cdot\frac1{(2\pi)^{N/2}}\exp\Big(-\frac12\big|A^{-1}(\underline y-\underline\mu)\big|^2\Big)\\
  &=\frac1{|\det A|}\cdot\frac1{(2\pi)^{N/2}}\exp\Big(-\frac12\big(A^{-1}(\underline y-\underline\mu)\big)^T\big(A^{-1}(\underline y-\underline\mu)\big)\Big)\\
  &=\frac1{|\det A|}\cdot\frac1{(2\pi)^{N/2}}\exp\Big(-\frac12(\underline y-\underline\mu)^T(A^{-1})^TA^{-1}(\underline y-\underline\mu)\Big)\\
  &\overset{(b)}{=}\frac1{(2\pi)^{N/2}\sqrt{\det C}}\exp\Big(-\frac12(\underline y-\underline\mu)^TC^{-1}(\underline y-\underline\mu)\Big)
\end{aligned}
\]
for every $\underline y\in\mathbb R^N$, where (a) holds because
$\underline X\sim\mathcal N_N(\underline 0,I_N)$, and (b) because
\begin{itemize}
  \item $C=A^2$, so $\det C=\det(A^2)=(\det A)^2$ and $|\det A|=\sqrt{\det C}$;
  \item $A$ is symmetric, so $(A^{-1})^TA^{-1}=(A^{-1})^2=(A^2)^{-1}=C^{-1}$.\qedhere
\end{itemize}
\end{proof}

\begin{remark}[Every positive semidefinite matrix is a Gaussian covariance]\label{rem:gv-existence}
In general, $C$ is symmetric and positive semidefinite. Then there still exists
$A\in\mathcal M_N(\mathbb R)$ symmetric and positive semidefinite (not necessarily invertible) such
that $A^2=C$.\footnote{The same construction: $A=Q\Lambda^{1/2}Q^T$, now with $\lambda_i\ge0$.} So
we can take $\underline X\sim\mathcal N_N(\underline 0,I_N)$ and define
\[
  \underline Y=A\underline X+\underline\mu\ \Longrightarrow\ \underline Y\sim\mathcal N_N(\underline\mu,C),
\]
since $AA^T=A^2=C$. \emph{It suffices to give $\underline\mu\in\mathbb R^N$ and a symmetric,
positive semidefinite $C\in\mathcal M_N(\mathbb R)$ to define a Gaussian vector.} Together with
\cref{lem:gv-covariance-psd}: the covariance matrices of random vectors with components in $L^2$
are exactly the symmetric positive semidefinite matrices.\footnote{This sentence is added to the
handwritten notes.}
\end{remark}

\begin{remark}[Degenerate Gaussian vectors]\label{rem:gv-degenerate}
One can prove that if $C$ is not invertible, then $\underline Y$ is not absolutely continuous, and
its law is concentrated on a proper affine subspace of $\mathbb R^N$, namely
$\underline\mu+\operatorname{Im}C$, which is negligible for Lebesgue measure.\footnote{The
handwritten notes say ``a proper subset of $\mathbb R^N$''. Every law on $\mathbb R^N$ is
concentrated on a proper subset (remove one point of probability zero); the point is that the
subset is a proper affine subspace, which has Lebesgue measure zero. Made precise here.} In this
case $\underline Y$ is called \emph{degenerate}.
\end{remark}

The notes state this without proof; \exref{ex:gv-degenerate} proves it, and shows that a degenerate
Gaussian vector is a non-degenerate one of dimension $\operatorname{rank}C$, placed in
$\mathbb R^N$ by an affine map. See also \citet[Section~29, ``Normal Distributions in
$R^k$'']{billingsley} and \citet[Section~3.10]{durrett}. Together with \cref{thm:gv-density}: a
Gaussian vector is absolutely continuous if and only if its covariance matrix is invertible. In
dimension $N=1$ we recover \textcircled{\scriptsize 4} of \cref{sec:gv-recap}.

\subsection*{Examples}

\begin{example}[From the covariance matrix to the density]\label{exa:gv-density-from-covariance}
Let
\[
  \underline X\sim\mathcal N_2\left(\begin{pmatrix}0\\0\end{pmatrix},\begin{pmatrix}2&1\\1&1\end{pmatrix}\right).
\]
Is it really a Gaussian vector? Does it have a density?

$C=\begin{pmatrix}2&1\\1&1\end{pmatrix}$ has $\det C=1$, so $C$ is invertible; it is symmetric and
positive definite.\footnote{Its two eigenvalues have product $\det C=1>0$ and sum
$\operatorname{tr}C=3>0$, so both are positive.} Hence $\underline X$ is a Gaussian vector
($C=A^2$ and $\underline X=A\underline Y+\underline 0$ with
$\underline Y\sim\mathcal N_2(\underline 0,I_2)$, \cref{rem:gv-existence}); moreover, it is
absolutely continuous (\cref{thm:gv-density}) and
\[
  f_{\underline X}(\underline x)=\frac1{2\pi\sqrt{\det C}}\exp\Big(-\frac12(\underline x-\underline 0)^TC^{-1}(\underline x-\underline 0)\Big)
  =\frac1{2\pi}\,e^{-\frac12\underline x^TC^{-1}\underline x}.
\]
Here
\[
  C^{-1}=\begin{pmatrix}1&-1\\-1&2\end{pmatrix},\qquad
  \underline x^TC^{-1}\underline x=(C^{-1})_{11}x_1^2+2(C^{-1})_{12}x_1x_2+(C^{-1})_{22}x_2^2
  =x_1^2-2x_1x_2+2x_2^2,
\]
so
\[
  f_{\underline X}(\underline x)=\frac1{2\pi}\,e^{-\frac12(x_1^2-2x_1x_2+2x_2^2)}\qquad
  \forall\underline x=(x_1,x_2)^T\in\mathbb R^2 .
\]
\end{example}

\begin{example}[From the density to the parameters]\label{exa:gv-covariance-from-density}
Let $\underline X$ be a random vector of $\mathbb R^2$ with density
\[
  f_{\underline X}(\underline x)=\frac1{2\pi\sqrt3}\,e^{-\frac13(x_1^2+x_1x_2+x_2^2)}\qquad
  \forall\underline x=(x_1,x_2)^T\in\mathbb R^2 .
\]
Is $\underline X\sim\mathcal N_2(\underline\mu,C)$, and for which $\underline\mu$ and $C$?

Comparing with \eqref{eq:gv-density} for $N=2$, we must have $\det C=3$ (compare the constants in
front of the exponentials, $\frac1{2\pi\sqrt{\det C}}=\frac1{2\pi\sqrt3}$), and
\[
  \frac12(\underline x-\underline\mu)^TC^{-1}(\underline x-\underline\mu)=\frac13(x_1^2+x_1x_2+x_2^2)
  \qquad\forall\underline x\in\mathbb R^2 .
\]
There are no terms of degree lower than $2$ on the right, so $\underline\mu=\underline 0$.%
\footnote{Expanding, the left-hand side is
$\frac12\underline x^TC^{-1}\underline x-\underline\mu^TC^{-1}\underline x+\frac12\underline\mu^TC^{-1}\underline\mu$.
The term of degree $1$, $-\underline\mu^TC^{-1}\underline x$, must vanish for every
$\underline x$, so $C^{-1}\underline\mu=\underline 0$ and $\underline\mu=\underline 0$. Added to
the handwritten notes.} Hence we must have
\[
  \underline x^TC^{-1}\underline x=\frac23(x_1^2+x_1x_2+x_2^2)
  =(C^{-1})_{11}x_1^2+2(C^{-1})_{12}x_1x_2+(C^{-1})_{22}x_2^2 ,
\]
that is,
\[
  C^{-1}=\begin{pmatrix}\frac23&\frac13\\[2pt]\frac13&\frac23\end{pmatrix}=\frac13\begin{pmatrix}2&1\\1&2\end{pmatrix}.
\]
Then $\det C^{-1}=\frac13$; $C^{-1}$ is invertible, and $C=(C^{-1})^{-1}$ has determinant $3$, as it
must. Hence
\[
  C=\begin{pmatrix}2&-1\\-1&2\end{pmatrix}\qquad\text{and}\qquad
  \underline X\sim\mathcal N_2\left(\begin{pmatrix}0\\0\end{pmatrix},\begin{pmatrix}2&-1\\-1&2\end{pmatrix}\right).
\]
(Conversely, this $C$ is symmetric and positive definite, and with $\underline\mu=\underline 0$ the
density \eqref{eq:gv-density} is exactly $f_{\underline X}$; two random vectors with the same
density have the same law, so $\underline X\sim\mathcal N_2(\underline 0,C)$ indeed.)\footnote{This
check is added to the handwritten notes.}
\end{example}

\begin{example}[A degenerate Gaussian vector]\label{exa:gv-degenerate}
Let $X_1,X_2$ be \iid with law $\mathcal N(0,1)$, and let
$\underline Y=A\underline X+\underline\mu$ with
\[
  \underline\mu=\begin{pmatrix}0\\0\end{pmatrix},\qquad A=\begin{pmatrix}1&1\\2&2\end{pmatrix}:\qquad
  \underline Y=\begin{pmatrix}1&1\\2&2\end{pmatrix}\begin{pmatrix}X_1\\X_2\end{pmatrix}
  =\begin{pmatrix}X_1+X_2\\2X_1+2X_2\end{pmatrix}=\begin{pmatrix}Y_1\\Y_2\end{pmatrix}.
\]
Then $\underline Y\sim\mathcal N_2(\underline 0,C)$ with
\[
  C=AA^T=\begin{pmatrix}1&1\\2&2\end{pmatrix}\begin{pmatrix}1&2\\1&2\end{pmatrix}=\begin{pmatrix}2&4\\4&8\end{pmatrix}.
\]
$\det C=16-16=0$: $C$ is not invertible. $\underline Y$ is degenerate, so it is not absolutely
continuous: there is no density $f_{\underline Y}=f_{(Y_1,Y_2)}$. Indeed, observe that $Y_2=2Y_1$:
\[
  1=\mathbb P(Y_2=2Y_1)=\mathbb P\big((Y_1,Y_2)\in D\big),\qquad D=\{(y_1,y_2)\in\mathbb R^2:y_2=2y_1\},
\]
and $D$ is a line! The law of $\underline Y$ is concentrated on the line $y_2=2y_1$.\footnote{This
line has Lebesgue measure zero in $\mathbb R^2$: its sections $\{y_2:(y_1,y_2)\in D\}=\{2y_1\}$ are
points, so $(m\otimes m)(D)=\int_{\mathbb R}m(\{2y_1\})\,\dd y_1=0$ by \cref{thm:pm-product-measure}.
Hence $\mathbb P_{\underline Y}$ is not absolutely continuous, as \cref{rem:gv-degenerate} says; the
line is $\underline\mu+\operatorname{Im}C$, the subspace spanned by $(1,2)^T$. Added to the
handwritten notes.}

\emph{However}, we shall see shortly (\cref{prop:gv-components}) that, reading the diagonal of $C$,
\[
  Y_1\sim\mathcal N(0,2),\qquad Y_2\sim\mathcal N(0,8):
\]
the densities $f_{Y_1}$ and $f_{Y_2}$ exist.
\end{example}

\begin{figure}[!htbp]
  \centering
  \includegraphics[width=\textwidth]{ch13-densities}
  \caption{The three examples. For the invertible covariance matrices of
    \cref{exa:gv-density-from-covariance,exa:gv-covariance-from-density}, a sample and three level
    sets of the density \eqref{eq:gv-density}: the ellipses
    $\{\underline y:\underline y^TC^{-1}\underline y\le r\}$ that hold $50\%$, $90\%$ and $99\%$ of
    the law (the quadratic form has the law $\chi^2(2)$, \exref{ex:gv-sampling}). For the singular
    $C$ of \cref{exa:gv-degenerate}, every sample point lies on the line $y_2=2y_1$.}
  \label{fig:gv-densities}
\end{figure}

\begin{remark}[computational]\label{rem:gv-code-sampling}
\coderef{code/src/probnotes/gaussian.py} samples $\mathcal N_N(\underline\mu,C)$ as in
\cref{def:gv-gaussian-vector}, as $A\underline X+\underline\mu$ with $AA^T=C$ and $\underline X$
standard. \texttt{sample\_gaussian} offers three choices of $A$: the symmetric square root of the
proof of \cref{thm:gv-density} (\texttt{psd\_sqrt}), the lower-triangular Cholesky factor
(\texttt{psd\_cholesky}), and $Q\Lambda^{1/2}$ from the spectral decomposition. All three also work
for singular $C$, where \texttt{numpy.linalg.cholesky} fails; they give the same law, but different
maps from $\underline X$ to $\underline Y$. \texttt{gaussian\_density} evaluates
\eqref{eq:gv-density}, and refuses a singular $C$; \texttt{support\_basis} and
\texttt{distance\_to\_support} describe the subspace $\underline\mu+\operatorname{Im}C$ of
\cref{rem:gv-degenerate}. See \exref{ex:gv-sampling} and \exref{ex:gv-whitening}.
\end{remark}

\demo{13-gaussian-vectors.ipynb}{a bivariate Gaussian vector with sliders for the variances and
the correlation, sampled through the symmetric square root and through the Cholesky factor, and
its image under an affine map, which becomes degenerate when the map is singular.}

\section{Components of a Gaussian vector and independence}\label{sec:gv-independence}

\begin{proposition}[The components are Gaussian]\label{prop:gv-components}
Let $\underline X=(X_1,\dots,X_N)^T\sim\mathcal N_N(\underline\mu,C)$. Then
\[
  X_j\sim\mathcal N(\mu_j,c_{jj})\qquad\forall j=1,\dots,N .
\]
\end{proposition}

\begin{proof}
Let $\Phi_{\underline X}$ be the characteristic function of $\underline X$:
\[
  \Phi_{\underline X}(\underline t)=\exp\Big(i\underline t^T\underline\mu-\frac12\underline t^TC\underline t\Big)
  \qquad\forall\underline t\in\mathbb R^N
\]
(\cref{prop:gv-characteristic}). Then, by \cref{prop:cf-marginals}, for every $t\in\mathbb R$
\[
  \Phi_{X_j}(t)=\Phi_{\underline X}(t\underline e_j)
  =\exp\Big(it\underbrace{\underline e_j^T\underline\mu}_{=\,\mu_j}-\frac12t^2\underbrace{\underline e_j^TC\underline e_j}_{=\,c_{jj}}\Big)
  =\exp\Big(it\mu_j-\frac12t^2c_{jj}\Big),
\]
which is the characteristic function of $\mathcal N(\mu_j,c_{jj})$. Since the characteristic
function characterises the law (\cref{thm:cf-uniqueness}), $X_j\sim\mathcal N(\mu_j,c_{jj})$.
\end{proof}

\begin{definition}[Jointly Gaussian random variables]\label{def:gv-jointly-gaussian}
Real random variables $X_1,\dots,X_N$ are \emph{jointly Gaussian} if
$\underline X=(X_1,\dots,X_N)^T$ is a Gaussian vector.
\end{definition}

By \cref{prop:gv-components}, jointly Gaussian random variables are Gaussian.

\begin{remark}\label{rem:gv-converse}
The converse does not hold: given $N$ Gaussian random variables $X_1,\dots,X_N$, they need not be
jointly Gaussian.
\end{remark}

\begin{example}[Gaussian marginals, not jointly Gaussian]\label{exa:gv-not-joint}
Let $X\sim\mathcal N(0,1)$ and
\[
  Y=ZX,\qquad\text{where }Z\sim U(\{-1,1\}),\ \text{that is, }\mathbb P(Z=1)=\mathbb P(Z=-1)=\tfrac12,
  \ \text{and }Z\indep X .
\]
Let $A\in\mathcal B(\mathbb R)$. Then, by the formula of total probability,
\[
\begin{aligned}
  \mathbb P_Y(A)&=\mathbb P(Y\in A)=\mathbb P(ZX\in A)\\
  &=\mathbb P(ZX\in A\mid Z=1)\,\mathbb P(Z=1)+\mathbb P(ZX\in A\mid Z=-1)\,\mathbb P(Z=-1)\\
  &=\mathbb P(X\in A\mid Z=1)\cdot\frac12+\mathbb P(-X\in A\mid Z=-1)\cdot\frac12
  \overset{(a)}{=}\frac12\big(\mathbb P(X\in A)+\mathbb P(-X\in A)\big)\\
  &\overset{(b)}{=}\mathbb P(X\in A)=\mathbb P_X(A),
\end{aligned}
\]
where (a) holds because $Z\indep X$, and (b) because $\mathbb P(-X\in A)=\mathbb P(X\in A)$, since
$-X\sim X$. Hence $\mathbb P_Y=\mathbb P_X$, that is,
\[
  Y\sim\mathcal N(0,1).
\]
The vector $(X,Y)^T$, however, is not a Gaussian vector. If it were, $X+Y$ would also be a
Gaussian random variable (\cref{thm:gv-affine}), since
\[
  X+Y=\begin{pmatrix}1&1\end{pmatrix}\begin{pmatrix}X\\Y\end{pmatrix}.
\]
But $X+Y$ is not Gaussian:
\[
  \mathbb P(X+Y=0)=\mathbb P(X+ZX=0)=\mathbb P\big(X(Z+1)=0\big)\overset{(c)}{=}\mathbb P(Z+1=0)=\frac12,
\]
where (c) holds because $\{X(Z+1)=0\}=\{X=0\}\cup\{Z=-1\}$ and $\mathbb P(X=0)=0$, $X$ being
absolutely continuous. So $X+Y$ is not absolutely continuous, and it is not almost surely constant
either, since $\mathbb P(X+Y=0)=\frac12<1$. Hence $X+Y$ is not Gaussian (by the last observation of
\cref{sec:gv-recap}),\footnote{The handwritten notes conclude from ``$X+Y$ is not absolutely
continuous'' alone. A degenerate Gaussian random variable is not absolutely continuous either, so
one must also exclude that $X+Y$ is almost surely constant. Corrected here.} and $(X,Y)^T$ is not a
Gaussian vector.
\end{example}

The law of $X+Y$ is in fact $\frac12\delta_0+\frac12\mathcal N(0,4)$, while a Gaussian vector with
the same marginals and the same covariance would give $X+Y\sim\mathcal N(0,2)$; see
\cref{fig:gv-not-joint}. The pair $(X,Y)$ is also an example of uncorrelated random variables that
are not independent (\exref{ex:gv-random-sign}); by the next theorem, this cannot happen for a
Gaussian vector.\footnote{This paragraph is added to the handwritten notes.}

\begin{figure}[!htbp]
  \centering
  \includegraphics[width=\textwidth]{ch13-not-joint}
  \caption{\Cref{exa:gv-not-joint}. Left: the pair $(X,ZX)$ lives on the two diagonals $y=\pm x$,
    although both coordinates are $\mathcal N(0,1)$. Right: the law of $X+ZX$ (simulated,
    $10^5$ draws), an atom of mass $\frac12$ at $0$ plus half of the $\mathcal N(0,4)$ density,
    against the $\mathcal N(0,2)$ density it would have if $(X,ZX)$ were a Gaussian vector.}
  \label{fig:gv-not-joint}
\end{figure}

\begin{theorem}[Independence of the components]\label{thm:gv-independence}
Let $\underline X=(X_1,\dots,X_N)^T\sim\mathcal N_N(\underline\mu,C)$. The components
$X_1,\dots,X_N$ are independent if and only if $C$ is diagonal, that is,
$c_{ij}=\operatorname{cov}(X_i,X_j)=0$ for all $i\ne j$.
\end{theorem}

\begin{proof}
``$\Rightarrow$'': if $X_1,\dots,X_N$ are independent, they are pairwise uncorrelated
(\cref{thm:pm-expectation-product}; the $X_i$ are in $L^2$ by \cref{rem:gv-mean-covariance}), so
$c_{ij}=0$ for $i\ne j$.

``$\Leftarrow$'': let $\Phi_{\underline X}$ be the characteristic function of $\underline X$:
\[
  \Phi_{\underline X}(\underline t)=\exp\Big(i\underline t^T\underline\mu-\frac12\underline t^TC\underline t\Big)
  =\exp\Big(i\sum_{j=1}^Nt_j\mu_j-\frac12\sum_{j,k=1}^Nc_{jk}t_jt_k\Big)
  \qquad\forall\underline t=(t_1,\dots,t_N)^T\in\mathbb R^N .
\]
By hypothesis $c_{jk}=0$ if $j\ne k$, so
\[
\begin{aligned}
  \Phi_{\underline X}(\underline t)&=\exp\Big(i\sum_{j=1}^Nt_j\mu_j-\frac12\sum_{j=1}^Nc_{jj}t_j^2\Big)
  =\exp\Big(\sum_{j=1}^N\Big(it_j\mu_j-\frac12c_{jj}t_j^2\Big)\Big)\\
  &=\prod_{j=1}^N\exp\Big(it_j\mu_j-\frac12c_{jj}t_j^2\Big)=\prod_{j=1}^N\Phi_{X_j}(t_j),
\end{aligned}
\]
by \cref{prop:gv-components}. By \cref{thm:cf-independence-criterion}, $X_1,\dots,X_N$ are
independent.
\end{proof}

The theorem also holds ``by blocks''. Suppose that $\underline X$ is split into two blocks,
\[
  \underline X=\begin{pmatrix}\underline X_1\\\underline X_2\end{pmatrix},\qquad
  \text{where $\underline X_1$ takes values in $\mathbb R^{N_1}$ and $\underline X_2$ in
  $\mathbb R^{N_2}$, with $N_1+N_2=N$.}
\]
Then
\[
  \underline\mu=\begin{pmatrix}\underline\mu_1\\\underline\mu_2\end{pmatrix},\qquad
  C=\begin{pmatrix}C_{11}&C_{12}\\C_{21}&C_{22}\end{pmatrix},
\]
where $\underline\mu_i$ is the mean vector of $\underline X_i$, $C_{ii}$ is the covariance matrix of
$\underline X_i$, and $C_{12}=C_{21}^T$ contains the covariances between the components of
$\underline X_1$ and those of $\underline X_2$. For example, for $N=3$, $N_1=2$ and $N_2=1$:
\[
  \underline X=\left(\begin{array}{c}X_1\\X_2\\\hline X_3\end{array}\right)\sim\mathcal N_3\left(
  \left(\begin{array}{c}\mu_1\\\mu_2\\\hline\mu_3\end{array}\right),
  \left(\begin{array}{cc|c}c_{11}&c_{12}&c_{13}\\c_{21}&c_{22}&c_{23}\\\hline c_{31}&c_{32}&c_{33}\end{array}\right)\right),
\]
with $\underline X_1=(X_1,X_2)^T$, $\underline X_2=X_3$, $\underline\mu_1=(\mu_1,\mu_2)^T$,
$\underline\mu_2=\mu_3$, $C_{11}=\left(\begin{smallmatrix}c_{11}&c_{12}\\c_{21}&c_{22}\end{smallmatrix}\right)$,
$C_{12}=\left(\begin{smallmatrix}c_{13}\\c_{23}\end{smallmatrix}\right)$,
$C_{21}=(c_{31}\ c_{32})$ and $C_{22}=c_{33}$.

\begin{corollary}[Independence by blocks]\label{cor:gv-blocks}
Let $\underline X=(\underline X_1^T,\underline X_2^T)^T\sim\mathcal N_N(\underline\mu,C)$, split
into blocks as above. Then
\[
  \underline X_1\indep\underline X_2\iff C_{12}=C_{21}^T=0 .
\]
\end{corollary}

\begin{proof}
Write $\underline t=(\underline t_1^T,\underline t_2^T)^T$ with $\underline t_i\in\mathbb R^{N_i}$.
Reasoning as before,
\[
\begin{aligned}
  \Phi_{\underline X}(\underline t)&=\exp\Big(i\sum_{j=1}^Nt_j\mu_j-\frac12\sum_{j,k=1}^Nc_{jk}t_jt_k\Big)\\
  &=\exp\Big(i\underline t_1^T\underline\mu_1+i\underline t_2^T\underline\mu_2
   -\frac12\underline t_1^TC_{11}\underline t_1-\frac12\underline t_2^TC_{22}\underline t_2
   -\frac12\underline t_1^TC_{12}\underline t_2-\frac12\underline t_2^TC_{21}\underline t_1\Big)\\
  &=\Phi_{\underline X_1}(\underline t_1)\,\Phi_{\underline X_2}(\underline t_2)\cdot
   \exp\Big(-\frac12\underline t_1^TC_{12}\underline t_2-\frac12\underline t_2^TC_{21}\underline t_1\Big),
\end{aligned}
\]
because $\underline X_i\sim\mathcal N_{N_i}(\underline\mu_i,C_{ii})$ (by \cref{thm:gv-affine}, with
$B=(I_{N_1}\ 0)$ or $B=(0\ I_{N_2})$), so that
$\Phi_{\underline X_i}(\underline t_i)=\exp\big(i\underline t_i^T\underline\mu_i-\frac12\underline t_i^TC_{ii}\underline t_i\big)$.
Hence
\[
  \Phi_{\underline X}(\underline t)=\Phi_{\underline X_1}(\underline t_1)\,\Phi_{\underline X_2}(\underline t_2)
  \quad\forall\underline t\in\mathbb R^N\iff C_{12}=C_{21}=0 .
\]
Indeed, $\underline t_2^TC_{21}\underline t_1=\underline t_1^TC_{12}\underline t_2$ (a real number
equals its transpose, and $C_{21}^T=C_{12}$), so the factorisation holds if and only if
$e^{-\underline t_1^TC_{12}\underline t_2}=1$, that is, $\underline t_1^TC_{12}\underline t_2=0$, for
all $\underline t_1,\underline t_2$; taking $\underline t_1,\underline t_2$ among the vectors of the
canonical bases, this means $C_{12}=0$.\footnote{This step is added to the handwritten notes.} By the
independence criterion (\cref{thm:cf-independence-criterion}, in its version for two blocks, which
has the same proof),\footnote{$\underline X_1\indep\underline X_2$ if and only if
$\Phi_{\underline X}(\underline t)=\Phi_{\underline X_1}(\underline t_1)\Phi_{\underline X_2}(\underline t_2)$
for every $\underline t$: replace $\mathbb P_{X_1}\otimes\dots\otimes\mathbb P_{X_N}$ by
$\mathbb P_{\underline X_1}\otimes\mathbb P_{\underline X_2}$ on
$\mathbb R^{N_1}\times\mathbb R^{N_2}$ in the proof of \cref{thm:cf-independence-criterion}. Added to
the handwritten notes.}
\[
  \underline X_1\indep\underline X_2\iff C_{12}=C_{21}=0 .\qedhere
\]
\end{proof}

From what we have just written we also obtain the following result.

\begin{proposition}[Concatenation of independent Gaussian vectors]\label{prop:gv-concatenation}
If $\underline X_1\sim\mathcal N_{N_1}(\underline\mu_1,C_{11})$,
$\underline X_2\sim\mathcal N_{N_2}(\underline\mu_2,C_{22})$ and
$\underline X_1\indep\underline X_2$, then
$\underline X=\left(\begin{smallmatrix}\underline X_1\\\underline X_2\end{smallmatrix}\right)$ is a
random vector with values in $\mathbb R^N$, $N=N_1+N_2$, and
\[
  \underline X\sim\mathcal N_N\left(\begin{pmatrix}\underline\mu_1\\\underline\mu_2\end{pmatrix},C\right),
  \qquad\text{where}\quad C=\begin{pmatrix}C_{11}&0\\0&C_{22}\end{pmatrix}.
\]
\end{proposition}

\begin{proof}
By independence,\footnote{As in the footnote to the proof of \cref{prop:cf-independence}: the
real and imaginary parts of $e^{i\underline t_1^T\underline X_1}$ and of
$e^{i\underline t_2^T\underline X_2}$ are bounded and measurable with respect to
$\sigma(\underline X_1)$ and $\sigma(\underline X_2)$. The notes leave this proof to the reader
(``from what we have just written''); added here.}
$\Phi_{\underline X}(\underline t)=\mathbb E\big[e^{i\underline t_1^T\underline X_1}e^{i\underline t_2^T\underline X_2}\big]
=\Phi_{\underline X_1}(\underline t_1)\,\Phi_{\underline X_2}(\underline t_2)$, and the computation in
the proof of \cref{cor:gv-blocks}, read with $C_{12}=C_{21}=0$, shows that this is
$\exp\big(i\underline t^T\underline\mu-\frac12\underline t^TC\underline t\big)$. The matrix $C$ is
symmetric and positive semidefinite, because $C_{11}$ and $C_{22}$ are
($\underline a^TC\underline a=\underline a_1^TC_{11}\underline a_1+\underline a_2^TC_{22}\underline a_2$),
so $\mathcal N_N(\underline\mu,C)$ exists (\cref{rem:gv-existence}) and has this characteristic
function (\cref{prop:gv-characteristic}). By \cref{thm:cf-uniqueness},
$\underline X\sim\mathcal N_N(\underline\mu,C)$.
\end{proof}

\begin{remark}[computational]\label{rem:gv-code-independence}
The module \coderef{code/src/probnotes/gaussian.py} samples the pair $(X,ZX)$ of
\cref{exa:gv-not-joint} (\texttt{sample\_random\_sign\_pair}), \texttt{random\_sign\_sum\_cdf} is the distribution function
of $X+ZX$, and \texttt{ks\_distance} computes $\sup_x|F_n(x)-F(x)|$ for an empirical distribution
function $F_n$, also when $F$ has atoms (\exref{ex:gv-random-sign}). \texttt{sample\_sign\_flip\_pair}
gives another pair of uncorrelated, dependent standard Gaussian variables
(\exref{ex:gv-sign-flip}). \texttt{cf\_factorisation\_defect} measures, on a sample of pairs, how far
the empirical characteristic function is from the product form of
\cref{thm:cf-independence-criterion}: it is of order $n^{-1/2}$ for the independent components of
a Gaussian vector with diagonal $C$, and stays near $\frac12$ for $(X,ZX)$. \texttt{is\_diagonal},
\texttt{blocks\_uncorrelated} and \texttt{concat\_independent} implement
\cref{thm:gv-independence,cor:gv-blocks,prop:gv-concatenation}.
\end{remark}

\demo{13-gaussian-vectors.ipynb}{pairs of standard Gaussian variables that are not jointly
Gaussian, with the law of their sum and a slider for the construction, and the whitening of a
Gaussian sample.}

\furtherreading
\textcite{billingsley}, Section~29, defines the normal distributions on $\mathbb R^k$ through their
characteristic functions (the notes construct them as $A\underline X+\underline\mu$ and
derive the characteristic function), and uses them for the multivariate central limit theorem
(\cref{ch:17-clt-delta}); \textcite{durrett}, Section~3.10, does the same. \textcite{ross-5},
Chapters~6 and~7, treats the bivariate normal law and the joint law of the sample mean and the
sample variance of a Gaussian sample at an elementary level (\exref{ex:gv-bivariate} and
\exref{ex:gv-sample-mean}).
